
Par
passage aux gradués associés dans (3.2.1) et (3.2.2) on obtient un homomorphisme de degré @@M0@@ :
@@M1@@
et grâce à (1.3.2) une formule de projection graduée
@@M2@@

\subsection*{3.3.}
Considérons enfin un morphisme @@M3@@ de Tor-dimension finie de schémas noethériens. Alors on sait (Exp. IV) que @@M4@@ induit un homomorphisme
@@M5@@
tel que
@@M6@@
pour tout @@M7@@-module cohérent @@M8@@.

Les @@M9@@-modules @@M10@@ sont concentrés sur @@M11@@ et lorsque les fibres de @@M12@@ sont partout de dimension @@M13@@, alors on sait (EGA IV 5.6.7) que
@@M14@@
Il s'ensuit que @@M15@@ induit pour tout @@M16@@ un homomorphisme
@@M17@@
d'où un homomorphisme de degré @@M18@@ de groupes abéliens gradués
